\documentclass[12pt]{article}
\usepackage{amsmath,amssymb,amsfonts}
\usepackage[margin=1in]{geometry}

\begin{document}

\textbf{Chapter 8, Problem 5.} Multiplying both sides of Eq.\ (8.27) by $V^{1/2}(r)$ [assume $V(r)$ to be strictly positive], show that if we define
\[
\tilde\phi(r) = V^{1/2}(r)\phi(r), \quad \tilde\phi_0(r) = V^{1/2}(r)\phi_0(r),
\]
and
\[
\tilde k(r,r') = V^{1/2}(r)G_0(r,r')V^{1/2}(r'),
\]
then
\[
\tilde\phi = \tilde\phi_0 + \int \tilde k(r,r')\tilde\phi(r')\,dr'.
\]
Calling $\tilde K$ the integral operator whose kernel is $\tilde k(r,r')$, show that the norm of $\tilde K$ is the same as the norm of the kernel of Eq.\ (8.29) obtained by the weight-function method. Also show that if we iterate the equation for $\tilde\phi(r)$, the object $V^{1/2}(r)$ never actually appears in the calculation---the ordinary Born series for $\phi(r)$ is duplicated, apart from a multiplicative factor of $V^{1/2}(r)$.

\bigskip
\textbf{Solution.}

The Lippmann--Schwinger equation (Eq.\ 8.27) is:
\[
\phi(r) = \phi_0(r) + \int G_0(r,r')V(r')\phi(r')\,dr'.
\]

\textbf{Step 1: Deriving the modified equation.}
Multiply both sides by $V^{1/2}(r)$:
\[
V^{1/2}(r)\phi(r) = V^{1/2}(r)\phi_0(r) + V^{1/2}(r)\int G_0(r,r')V(r')\phi(r')\,dr'.
\]
Write $V(r') = V^{1/2}(r')V^{1/2}(r')$ and insert into the integral:
\[
\tilde\phi(r) = \tilde\phi_0(r) + \int V^{1/2}(r)G_0(r,r')V^{1/2}(r')\cdot V^{1/2}(r')\phi(r')\,dr'.
\]
Since $V^{1/2}(r')\phi(r') = \tilde\phi(r')$:
\[
\boxed{\tilde\phi(r) = \tilde\phi_0(r) + \int \tilde k(r,r')\tilde\phi(r')\,dr'.} \quad\checkmark
\]

\textbf{Step 2: Norm equivalence.}
The Hilbert--Schmidt norm of $\tilde K$ is:
\[
\|\tilde K\|^2 = \int\!\!\int |\tilde k(r,r')|^2\,dr\,dr' = \int\!\!\int V(r)|G_0(r,r')|^2 V(r')\,dr\,dr'.
\]
The weight-function method (Eq.\ 8.29) uses the kernel $k_w(r,r') = G_0(r,r')V(r')$ with the weight function $w(r)=V(r)$, giving the weighted norm:
\[
\|K_w\|^2 = \int\!\!\int |G_0(r,r')|^2 V(r')^2\, V(r)\,dr\,dr'?
\]
Actually, in the weight-function method, the inner product is $\langle f,g\rangle = \int f^*(r)g(r)V(r)\,dr$, so the norm of the operator $K$ with kernel $G_0(r,r')V(r')$ in this weighted space is:
\[
\|K\|_w^2 = \int\!\!\int |G_0(r,r')|^2 V(r') V(r)\,dr\,dr',
\]
which is exactly $\|\tilde K\|^2$. $\checkmark$

\textbf{Step 3: Iteration reproduces the Born series.}
Iterating $\tilde\phi = \tilde\phi_0 + \tilde K\tilde\phi$:
\begin{align*}
\tilde\phi &= \tilde\phi_0 + \tilde K\tilde\phi_0 + \tilde K^2\tilde\phi_0 + \cdots
\end{align*}
The $n$-th term is:
\[
\tilde K^n\tilde\phi_0(r) = \int\!\cdots\!\int V^{1/2}(r)G_0(r,r_1)V(r_1)G_0(r_1,r_2)V(r_2)\cdots G_0(r_{n-1},r_n)V^{1/2}(r_n)\phi_0(r_n)\,dr_1\cdots dr_n.
\]
Since $V^{1/2}(r_n)\phi_0(r_n) = \tilde\phi_0(r_n)$, and using $V(r_i) = V^{1/2}(r_i)\cdot V^{1/2}(r_i)$, the factors of $V^{1/2}$ pair up internally, yielding:
\[
\tilde\phi^{(n)}(r) = V^{1/2}(r)\int\!\cdots\!\int G_0(r,r_1)V(r_1)G_0(r_1,r_2)V(r_2)\cdots G_0(r_{n-1},r_n)V(r_n)\phi_0(r_n)\,dr_1\cdots dr_n.
\]
This is exactly $V^{1/2}(r)$ times the $n$-th Born approximation term for $\phi(r)$. Thus the iteration of the $\tilde\phi$ equation reproduces the standard Born series for $\phi$ multiplied by $V^{1/2}(r)$, and $V^{1/2}$ never appears internally in the calculation. $\blacksquare$

\end{document}
